\clearpage
\section{指标与行列式线}\label{sec:11}\label{section:11}
\subsection{行列式线}\label{sec:11a}
本节的讨论主要围绕行列式线展开。行列式线是对应于 Fredholm 算子的一维实向量空间。实际上，我们关心的是这条线的两种可能定向所成的集合，而非线本身。特别地，每当我们说两个行列式线（或稍后引入的类似对象）之间的同构是典范的，总是指在相差一个正实数倍的意义下典范。作了这一措辞上的说明之后，设 $D$ 为两个实 Banach 空间之间的 Fredholm 算子。$D$ 的行列式线定义为其核与余核的最高次外幂的张量积，确切地说，
\begin{definition*}
\[
\det(D)=\lambda^{\mathrm{top}}(\operatorname{coker}(D)^\vee)\otimes\lambda^{\mathrm{top}}(\ker(D)).
\]
\end{definition*}
\SourcePage{158}
由于只关心定向，对余核取对偶是多余的；但我们仍愿意保留这种写法，使公式与复行列式线的公式相对应。\znote{150}{行列式线是 Fredholm 指标理论的标准工具。对 Fredholm 算子 $D:E\to F$，取 $\lambda(\ker D)\otimes\lambda(\operatorname{coker}D)^\vee$（$\lambda$ 表示最高次外幂）为它的行列式线；当算子连续形变时，这些一维空间拼成底空间上的线丛，即行列式线丛。nLab 把这个构造从有限维的行列式 $\det T:\wedge^nV\to\wedge^nW$ 出发加以推广，指出向量丛的最高次外幂丛正是行列式线丛，并由此讨论 Quillen 行列式线丛以及 Grassmann 流形上的行列式丛，可用作理解本条定义的背景。[来源：nLab "determinant line bundle"，https:/\allowbreak /\allowbreak ncatlab.org/\allowbreak nlab/\allowbreak show/\allowbreak determinant+\allowbreak line+\allowbreak bundle（2026-11-15 访问）]}
存在典范同构
\begin{equation}\label{eq:11.1}
\det(D_1)\otimes\det(D_2)\cong\det(D_1\oplus D_2),
\end{equation}
其定义是将所有向量视为奇向量，并使用 Koszul 符号规则。具体地，若 $\{v_{k,\allowbreak i}\}$ 是 $\ker(D_k)$ 的基，$\{w_{k,\allowbreak i}^\vee\}$ 是 $\operatorname{coker}(D_k)^\vee$ 的基，则同构为
\begin{equation}\label{eq:11.2}
\begin{aligned}
&\left(\bigwedge_iw_{1,i}^\vee\otimes\bigwedge_iv_{1,i}\right)
\otimes\left(\bigwedge_iw_{2,i}^\vee\otimes\bigwedge_iv_{2,i}\right)\\
&\quad\longmapsto(-1)^{\dim(\operatorname{coker}D_2)\cdot\operatorname{index}(D_1)}
\bigwedge_iw_{2,i}^\vee\wedge\bigwedge_iw_{1,i}^\vee
\otimes\bigwedge_iv_{1,i}\wedge\bigwedge_iv_{2,i}.
\end{aligned}
\end{equation}
它严格满足结合律，并在分次意义下满足交换律。后者指的是，两个同构的复合
\[
\det(D_1)\otimes\det(D_2)\cong\det(D_1\oplus D_2)
=\det(D_2\oplus D_1)\cong\det(D_2)\otimes\det(D_1)
\]
与直接交换因子的映射相差符号 $(-1)^{\operatorname{index}(D_1)\operatorname{index}(D_2)}$。

暂且考虑有限维向量空间之间的线性映射 $D:H\to K$。可以用唯一的方式选择识别 $t_D:\det(D)\to\det(0)=\lambda^{\mathrm{top}}(K^\vee)\otimes\lambda^{\mathrm{top}}(H)$，使之与 $D$ 在左、右两侧复合同构相容，并满足以下三条额外规则：若 $D=0$，则 $t_D$ 为恒等映射；若 $D=\id:H\to H$，其中 $H$ 为一维向量空间，则 $t_D:\R=\det(\id)\to H^\vee\otimes H$ 将 $1$ 映为 $v^\vee\otimes v$，这里 $\langle v^\vee,\allowbreak v\rangle>0$；最后，这些同构与直和相容，也就是说下图交换：
\SourceDiagram{diagram-11a}{0.46523027}
其中 $0_k$ 是与 $D_k$ 具有相同定义域和目标空间的零映射。另一种做法，是对任意 $D:H\to K$ 直接定义 $t_D$：取 $H$ 的基 $e_1,\allowbreak \ldots,\allowbreak e_m$ 与 $K$ 的基 $f_1,\allowbreak \ldots,\allowbreak f_n$，满足 $D(e_1)=f_1,\allowbreak \ldots,\allowbreak D(e_k)=f_k$、$D(e_{k+1})=0,\allowbreak \ldots,\allowbreak D(e_m)=0$，再取它们的对偶基 $e_i^\vee,\allowbreak f_i^\vee$。则
\begin{equation}\label{eq:11.3}
\begin{aligned}
t_D:\ &(f_n^\vee\wedge\cdots\wedge f_{k+1}^\vee)\otimes(e_{k+1}\wedge\cdots\wedge e_m)\in\det(D)\\
&\longmapsto(f_n^\vee\wedge\cdots\wedge f_1^\vee)\otimes(e_1\wedge\cdots\wedge e_m)\in\det(0).
\end{aligned}
\end{equation}
有了这些准备，回到一般情形。设 $\mathfrak F(H,\allowbreak K)$ 是实 Banach 空间 $H,\allowbreak K$ 之间所有 Fredholm 算子所成的空间。\emph{行列式线丛} $\underline{\det}\to\mathfrak F(H,\allowbreak K)$ 的纤维为 $\det(D)$，其拓扑结构可由下述性质刻画。首先，若 $H,\allowbreak K$ 有限维，同构 $t_D$ 应给出 $\underline{\det}$ 的一个平凡化。其次，上面定义的映射
\SourcePage{159}
$\det(D_1)\otimes\det(D_2)\cong\det(D_1\oplus D_2)$，是在 $\mathfrak F(H_1,\allowbreak K_1)\times\mathfrak F(H_2,\allowbreak K_2)$ 上的行列式线丛同构。最后，若限制到满射算子的子空间，则等式 $\det(D)=\lambda^{\mathrm{top}}(\ker(D))$ 给出行列式线丛与核向量丛的最高次外幂之间的同构；这是常用的通过稳定化定义行列式线丛的方法的一个小变体。不必固定两个 Banach 空间，也可以考虑两个 Banach 向量丛以及它们之间的 Fredholm 映射，从而得到底空间上的行列式线丛。
\begin{note}\label{st:11.1}
这里引入的符号约定采用\bcite{87}的约定；亦见\bcite{77}附录 D。在\bcite{36}第一卷第 62、362 页中，它被称为“观点 II”。提一下另一种可能的约定，也许有助于理解：将\eref{11.2}右侧换成（即去掉符号）
\[
\bigwedge_iw_{1,i}^\vee\wedge\bigwedge_iw_{2,i}^\vee\otimes\bigwedge_iv_{1,i}\wedge\bigwedge_iv_{2,i}
\]
并将\eref{11.3}中 $f_i^\vee$ 的次序反转。这样得到另一个线丛 $\underline{\det}'$，它的纤维与 $\underline{\det}$ 相同；两个线丛同构，却并非以平凡的方式同构，恒等映射 $\underline{\det}\to\underline{\det}'$ 并不连续。
\end{note}
